\chapter{Words Are Not Fixed Labels on Things}
\section{A Dictionary and the Word Da}
First pick up a dictionary. Any kind will do; a thicker one is better.

Turn to the Chinese word dǎ. You will find a long list of meanings that the editors have doggedly assembled: dǎ rén, hit someone (strike); dǎ shuǐ, fetch water (scoop); dǎ chē, take a taxi (hire); dǎ zhēn, get a shot (inject); dǎ qiú, play ball (play); dǎ gōng, work (do a job); dǎ jiāodao, have dealings (interact); dǎ bǐfang, draw an analogy (give an example); dǎ kēshuì, doze (grow sleepy); dǎ guānsi, bring a lawsuit (litigate); dǎ zhékòu, give a discount (reduce). More than twenty senses, each with examples, like twenty-odd ID photographs of the same character, every expression different.

Now ask yourself: have you ever looked up dǎ? If you grew up speaking Chinese, probably not. You could say ``get a shot'' when you were little and ``play games'' by school age. You would never understand ``I'm going downstairs to fetch water'' as ``I'm going downstairs to beat up water.'' Strangely, what a dictionary struggles to explain in twenty entries, your brain handles instantly: glance at the words surrounding dǎ, and you know which dǎ this is.

What sort of dictionary is inside your head? Who edited it?

Pick up another small object for an experiment: a sticky note. Write ``chair'' on it and slap it onto a chair's back. The act seems perfectly natural. Words are labels, after all: one for each thing, meaning whatever it is stuck to. Adults teach babies around the world this way: point at a lamp and say ``lamp,'' point at a dog and say ``dog.'' Words are names attached to things---our simplest picture of language.

This chapter will dismantle that picture. We will see that words are not labels stuck on things, or at least not labels fixed for life. They slide, move, get peeled off and attached elsewhere; different people attach different meanings. Some things carry dozens of overlapping labels. For others, you circle them with a label in hand and find you have no idea where to stick it.

\section{New York Harbor, 1886: A Tomato Goes to Court}
Now enter a scene.

Time: spring 1886. Place: New York Harbor, United States. Wooden crates crowd the wharf; dockworkers shout, and the salty sea air mingles with the sweet-sour smell of slowly fermenting produce. A cargo ship from the West Indies is unloading. Inside its crates are tomatoes: red, yellow, perfectly round.

The importers are the Nix family, experienced produce merchants. Customs official Hedden stands beside the crates with a tax bill: these goods are vegetables, subject to the prescribed ten-percent duty.

The basis is a tariff law Congress passed in 1883: imported vegetables are taxed; imported fruit is exempt. The Nixes object. Tomatoes are fruit! Bite one: it is juicy and contains seeds. Botanically, a fruit is the seed-bearing part developed from a flower's ovary, and a tomato fits perfectly. An apple is a fruit; so is a tomato. Why exempt apples but tax tomatoes?

Customs will not budge: tomatoes are vegetables. Unable to agree, the Nixes sue. The case climbs court after court, all the way to the United States Supreme Court in Washington. In 1893, the Court hears it.

What follows is an unusually peculiar scene in legal history. The principal evidence the lawyers set before the justices consists not of botanical specimens or laboratory reports, but dictionaries. One dictionary after another is opened in court; lawyers read aloud the definitions of ``fruit'' and ``vegetable,'' as though the dispute concerns the dictionary's authority rather than tomatoes.

The ruling arrives: the justices unanimously decide that tomatoes are vegetables. Justice Gray writes the opinion. Its reasoning, in essence, is that botanical classification is one thing, ordinary life another. In kitchens and at dining tables, tomatoes are eaten with the main meal, like peas, cucumbers, eggplants, and cabbages, not afterward as dessert like fruit. Thus, in the world governed by law, words take their everyday meanings.

A small episode in the hearing is worth remembering. Both sides call witnesses: seasoned dealers with decades in the produce trade. The justices ask not about botany but something simpler: in your trade's language, which category do tomatoes belong to? The dealers' answers are almost unanimous: we have always sold them as vegetables. Notice the question's form. Throughout, the court asks not ``What are tomatoes really?'' but ``What have people always called them?'' A place apparently devoted to ``definitions'' ends up relying on usage.

In short: in a botany textbook, a tomato is a fruit; on a customs form, a vegetable. Both answers are ``right,'' because they follow different sets of rules.

Do not rush to laugh at the Americans for having nothing better to do. The Nixes spent more on the lawsuit than the tax itself. What they were really contesting is this chapter's question:

\section{Who Decides What a Word Means?}
First make the conflict clear, without rushing to answer it.

Three parties stand in the tomato case, each convinced it holds the correct answer to ``What is a tomato?''

The botanist says: a word's meaning lives in the thing itself. Fruit has a strict definition; a tomato meets it, so a tomato is a fruit. This is a question of fact, not a vote. If most people say the moon is cheese, the moon will not become cheese.

The cook says: a word's meaning lives in its use. How a tomato grows matters less than what people do with it: chop it into salad, stew it in soup, scramble it with eggs. Language serves life, not microscopes.

Customs says: a word's meaning lives in the statute. Tariff law needs clear categories; collectors cannot convene a philosophy seminar every time they levy a duty.

Behind all three stands one real question: \textbf{where does a word's meaning live?} In a thing's nature, in the dictionary's authority, or in millions of people's daily usage?

The question runs deeper than it appears, because words inhabit a messier world than we imagine. Consider some other forms of disorder:

\textbf{One word, many meanings.} Take our opening dǎ: one character, more than twenty meanings, their connections sometimes barely visible. How did so many meanings grow? Are they still relatives?

\textbf{Synonyms, different lives.} Mǔqīn, māma, niáng---formal ``mother,'' familiar ``Mom,'' traditional ``Ma''---refer to the same person. Yet write ``my honored maternal lady'' in a school essay and your teacher will underline it; call out ``Comrade Mother'' at home and your mom will think you have got into trouble. Words with the ``same'' meaning wear quite different clothes and attend different occasions. Synonyms have never been truly identical: they share a referent but carry their own scents.

\textbf{Opposites have middle ground.} The opposite of ``cold'' is ``hot,'' yet lukewarm water is neither. The opposite of ``success'' is ``failure,'' yet most people's days are mostly neither. Antonyms are less like opposite sides of a wall than opposite slopes of a mountain, with a broad gray plain between them.

\textbf{Fuzzy boundaries.} How tall counts as ``tall''? Does 1.75 meters qualify? How many grains of rice make ``a heap''? Remove one grain: is it still a heap? No dictionary dares give you a number, because the word itself comes with none.

Now take a side. If you were the judge in the tomato case, how would you rule? Who should decide meaning: experts, dictionaries, or everyone who uses the word? Carry your answer onward and see how many sections it survives.

\section{Who Governs the Dictionary? Both Sides' Full Case}
This section offers two positions. We will reason seriously through each, including the one you may disagree with.

\textbf{Position one: somebody must regulate meaning.}

Let us give this side its full case. ``Language police'' may sound like a joke, but its reasons are no joke at all.

First, unstable meanings can cause real trouble in some settings. A contract says ``delivery'': exactly when has delivery occurred? Medicine instructions say ``a small amount'': how much? Legal documents, medical reports, and examination questions all assume that words have stable, shared meanings. Without that stability, a signature loses its point.

Second, if meaning is wholly free, bad actors always exploit the loopholes first. A seller defines ``sugar-free'' as ``no added sucrose,'' then adds high-fructose corn syrup; ``entirely handmade'' becomes ``a hand touched the machine's button.'' Let anyone stretch meanings at will, and language becomes modeling clay in a fraudster's hands.

Third, shared language is infrastructure for cooperation among hundreds of millions, like railway gauges. Track workers cannot individually decide how wide to lay the rails; speakers cannot entirely decide for themselves what words mean. Otherwise your bridge and mine will not line up, and the bridge will collapse.

This is not merely theory. France built an institution for it: the Académie française, founded by Cardinal Richelieu in 1635. Forty members known as ``the immortals'' undertake one task: compile the most authoritative French dictionary and settle disputes over French. They still meet today, promoting courriel, for example, and urging people not to use the English loan email. Smaller Iceland goes further. Icelandic originally had no word for ``computer,'' so a new one was forged from native roots: tölva, combining ``number'' and ``female prophet'' to mean ``a prophetess who calculates.'' In this view, language is an inheritance that must not be left to wash away under foreign words.

\textbf{Position two: language belongs to its users, and nobody can lock it down.}

This side has an equally complete case.

First, dictionaries record; they do not command. Editors run behind the crowd, writing down meanings people have already developed. Chinese gěilì, ``awesome'' or ``effective,'' became popular first and entered dictionaries afterward, never the other way round.

Second, English is the strongest evidence. The English-speaking world has never had an ``English Academy,'' no institution authorized to rule on English. Yet English became one of the planet's most widely used languages. Language maintains order not through police but through self-correction in hundreds of millions of daily uses by its speakers.

Third, it cannot be controlled: semantic drift is language breathing. Consider nice, ultimately from Latin nescius, ``ignorant.'' In English it first meant ``foolish,'' then ``fussy'' and ``fine,'' before reaching today's ``pleasant.'' Across centuries, the same word moved from insult to compliment without any parliament voting on it. In Chinese, xiǎojiě, ``Miss,'' was a respectable address decades ago but later acquired other associations in some settings; tóngzhì, ``comrade,'' has also shifted several times. You may dislike these changes, but you cannot stop them any more than you can stop a river changing course.

Who is right? First notice the power question hidden behind the dispute: \textbf{if someone must regulate meaning, why should it be those people?} The Académie française's forty seats hold writers, professors, and public figures, not market-stall speakers or young people inventing words. Yet ``nobody in charge'' does not mean equality either. A new meaning spreads when influential media, celebrities, and bloggers use it first. Ten thousand words invented by ordinary people may not outweigh one trending topic. There is no clean answer, only a gap we must watch closely: \textbf{meaning belongs to everyone, but powerful people have always reached first for the pen that defines it.}

\section{Thinking Bridge: Take a Word Apart}
Now learn a tool you can carry away. Next time you meet an unfamiliar word---in Chinese, English, or any language---do not panic. Try taking it apart.

Linguists call ``the smallest meaningful unit of language'' a \textbf{morpheme}. ``Smallest'' means that further division destroys the meaning. Chinese rén, ``person,'' is one morpheme: divide its written character into its two slanting strokes, and you have marks, not meanings. Bōli, ``glass,'' is also one morpheme: bō and li separately mean nothing here; the two characters form a meaningful partnership only together.

Morphemes play two roles. The \textbf{root} is a word's backbone, carrying its main meaning; an \textbf{affix} is a small attached part that modifies or reshapes it. Consider several groups:

Chinese: kě'ài, ``lovable''---ài, ``love,'' is the root; kě is a prefix meaning ``worthy of,'' so literally ``worthy of love.'' The zi in zhuōzi, ``table,'' and tou in shítou, ``stone,'' are suffixes with little meaning of their own, padding one-syllable characters into two-syllable words. The huà in xiàndàihuà, ``modernization,'' is a suffix meaning ``become a certain way.'' Thus lǜhuà means making the surroundings green; chǒuhuà means making an image ugly. The same part does the same job on different roots.

English: unhappy combines un-, ``not,'' with happy. The part tele- comes from Greek and means ``far'': telephone is ``sound from afar,'' television ``seeing afar,'' telescope an instrument for ``looking afar.'' Learn these pieces, and dozens of unfamiliar words surrender without a fight.

Arabic uses another assembly method. The root is three consonants, like a skeleton; insert different vowel patterns and you produce a word family. Take the writing-related skeleton k-t-b: kitab is ``book,'' kataba ``he wrote,'' maktab ``a place for writing'' or ``office,'' maktaba ``library,'' katib ``scribe.'' Press the same dough into different molds and you get a tray of pastries. Recognize the molds, and you recognize the whole tray.

Chinese most commonly forms words by \textbf{compounding}: joining two or more morphemes directly. Huǒchē, ``fire-vehicle'' (train); diànnǎo, ``electric-brain'' (computer); bīngxiāng, ``ice-box'' (refrigerator); xǐyījī, ``wash-clothes-machine'' (washing machine). The first time you hear xǐyījī, you need no dictionary: its parts assemble the meaning themselves. This is one of Chinese's superpowers: endlessly building new words from a limited stock of parts. Borrowed words often use \textbf{phonetic transcription}: shāfā (sofa), kāfēi (coffee), qiǎokèlì (chocolate), luóji (logic). The sound is borrowed, with the meaning carried over as a package.

But this tool has limits. Here is an easily misjudged counterexample: \textbf{not every word comes apart, and the meanings of the parts do not always add up to the whole.} Húdié, ``butterfly,'' and pútao, ``grape,'' cannot be split meaningfully: division leaves fragments. Shāfā certainly does not mean ``daydreaming on sand,'' despite what its individual characters might suggest; characters in phonetic loans are only shells for sound. Be especially wary of dōngxi, ``things'': its meaning has almost nothing to do with the directions ``east'' and ``west.'' Mǎhu, ``careless,'' likewise has nothing to do with horses and tigers. Parts sometimes add up to a word's meaning and sometimes do not. Use word analysis properly: take it apart to guess, then put it back into the sentence to test the guess, not to declare a verdict.

Here is a little exercise to take away. Next time you see these words, separate their parts before guessing: fǎnyīngxióng, ``anti-hero,'' a protagonist unlike a traditional hero; wěiqiúmí, ``pseudo-fan,'' someone not really a sports fan; kěhuíshōu, ``recyclable,'' worthy of or able to be recycled. Then try English rewrite (write again), preview (see beforehand), impossible (not possible). The parts inventory is small: only a few dozen common affixes, appearing over and over like connecting studs on building blocks. Recognize one new part each day, and your vocabulary starts multiplying itself.

\section{A Little Evidence from 2,300 Years Ago}
We possess a very old primary text for the dispute over ``where a word's meaning lives.''

Xunzi, a thinker of the late Warring States period, wrote ``Rectification of Names,'' specifically examining the relationship between names and things.

First, some background: why did he care? In his era the old order was falling apart and new doctrines were everywhere. ``Names'' and ``realities''---terms and what they designated---failed to line up: some held real power without matching titles; some things had changed while their names stayed old. A group of disputers specializing in names and realities was later called the School of Names. One, Gongsun Long, left the famous proposition ``A white horse is not a horse.'' His reasoning, roughly, was that ``horse'' denotes shape, ``white'' color; ``white horse'' combines color and shape and therefore differs from ``horse,'' which denotes shape alone. People have argued for over two thousand years whether this is sophistry or insight. It proves at least that Chinese thinkers were already seriously examining the relation between words and what they name. Xunzi's essay is one of that great debate's most mature answers. It includes this sentence:

\begin{quote}
Names have no inherent appropriateness: they are assigned by agreement. When agreement becomes established custom, they are called appropriate; what departs from agreement is called inappropriate. ---Xunzi, ``Rectification of Names''
\end{quote}
In essence: no name is ``suitable from the start.'' People agree to call something by it; once agreement and habit establish it, it becomes suitable. What violates that agreement is unsuitable.

Chew this sentence thoroughly. It has two layers, neither dispensable.

First: no natural bond joins words to things. A tomato is not called a tomato because the word is carved into it. A dog may be called gǒu, dog, or inu; names in every language arise from agreement. This layer directly sentences the ``label view'' to death: no label grows naturally on a thing.

The second layer is equally important and often missed: \textbf{once agreement is established, it binds.} Xunzi did not say ``so call anything whatever you like.'' Quite the opposite: ``what departs from agreement is called inappropriate.'' Once custom is established, insist on calling a dog a ``cat,'' and the mistake is yours, not the dog's. Meaning is society's shared property; individuals have no right to privatize it.

Return now to the US Supreme Court in 1893. The justices' reasoning effectively applies Xunzi's principle: in the world of tariff law and produce markets, the convention for ``vegetable'' belongs to the kitchen, not botany. A Chinese thinker more than two millennia ago and American justices more than a century ago, separated by half the globe and two civilizations, touched the same stone: \textbf{a word's meaning comes from a community's shared use; once that shared use forms, it becomes a real rule.}

\section{Why the Tomato Lost: Three Explanations}
One fact---the Supreme Court classified the tomato as a vegetable---can enter three very different explanations. First distinguish four things: what happened (the ruling, an undisputed historical fact); why it happened (explanation, where the three accounts below compete); how people then understood it (their own account, given in the opinion); and how we evaluate it (judgment: was it fair?). We are comparing the second layer.

\textbf{Explanation one: the court bowed to ignorance.} By scientific standards, the botanical definition is the ``true'' one. A tomato is a fruit; the ruling let the majority's mistaken usage override fact. This account's strength is its clear commitment: some things do not depend on what people call them. Earth does not flatten because people call it flat. Its limitation is equally clear: it cannot answer a practical question. Tariff law was never intended to classify plants, but to collect money according to how food is used. In the kitchen, tomatoes belong with peas, not apples. Collect customs duties through a microscope, and every penny will sit awkwardly.

\textbf{Explanation two: meaning belongs to use, and the court was right.} On the Xunzi-and-usage view, ``vegetable'' has always been a kitchen word, a word of daily life, not a scientific term. Its meaning comes from centuries of shared use. Tomatoes are used exactly like vegetables in ordinary life, so the ruling respected language's real rules. This account's strength is respecting how language actually works. Its vulnerability is this: if majority usage sets the standard, what happens when most people say ``whales are fish'' or ``bats are birds''? Its answer must be that everyday language and science have their own conventions and territories, and neither should police the other's ground. ``A whale is not a fish'' is a biology-class fact; displaying whale alongside fish in a fish market is not thereby fraud.

\textbf{Explanation three: the issue is tax, not words.} This view changes the angle: never mind the dictionaries the justices discuss; look at who stands behind the ruling. Classify tomatoes as fruit, and government loses that tariff revenue. More importantly, once tomatoes open the door to ``reclassification by scientific definition,'' cucumbers, eggplants, peppers, pea pods---a host of botanical fruits treated as vegetables at table---will all invite fresh lawsuits. Half the customs classification system could collapse. The court was persuaded by order and revenue, not dictionaries. This sharp account can explain questions of timing that the first two cannot. But it can go too far. The opinion is earnestly reasoned; the justices appear genuinely to believe that legal language should respect everyday language. Translate every motive into interest, and you lose sight of real conviction: a new misreading.

Each explanation holds a piece of truth. This is not a compromise for its own sake, but a methodological reminder: \textbf{when an explanation claims to explain everything, it has usually discarded something in advance.}

\section{Further Challenge: Why Is ``Game'' So Hard to Define?}
Now bring out the chapter's deepest theory. It addresses an irritating puzzle: why can humans not define some of their most ordinary words?

Try the traditional method first. Since Aristotle, Western logic has offered a standard format for definition: identify the larger class a word belongs to, then add the feature distinguishing it from others in that class. ``A human is a rational animal'': animal is the larger class, rational the distinguishing feature. Under this method, a perfect definition gives \textbf{necessary and sufficient conditions}. Whatever satisfies them is inside the circle; whatever does not is outside, cleanly divided.

Now try defining ``chair'' yourself. ``Furniture with four legs and a back, for sitting''? A beanbag has no legs: is it a chair? You can sit on a step: is a step a chair? Someone sits daily on an upturned wooden crate: is that a chair? Every condition you write invites a real-world exception. Plug each exception, and the definition grows more tangled, until even you no longer want to read it.

In the mid-twentieth century, Austrian philosopher Wittgenstein performed the same experiment with ``game'' in Philosophical Investigations, published in 1953, to even greater effect. Chess, cards, football, hide-and-seek, children spinning in circles, someone tossing a ball alone: find a feature shared by all these activities and unique to games. ``All have winners and losers''? Spinning does not. ``All have rules''? Children fooling around need not. ``All are for fun''? Professional chess players earn their living under enormous pressure. Search as you will, no feature runs through every game. Wittgenstein said the members of the ``game'' family resemble a large human family: the eldest and Grandma share crescent-shaped eyes; Grandma and a cousin share quick tempers; the cousin and the eldest walk alike. Members share resemblances, but no single feature belongs to the whole family. He called this \textbf{family resemblance}.

Around the same period, psychologist Eleanor Rosch equipped this idea with experimental data. Asked whether something was a bird, participants judged ``a robin is a bird'' quickly and confidently, ``a penguin is a bird'' more slowly and hesitantly, and ``a chicken is a bird'' more slowly still. This suggests that mental categories are less like sharp boundary lines than circles with centers and margins: \textbf{the most typical member, the prototype, stands at the center; members grow less typical toward the edge, but nobody can say exactly where the circumference lies.}

Color is the best example. Physically, a rainbow is a continuous spectrum, its wavelengths changing smoothly from long to short, with no dividing lines printed on it. Where is the boundary between ``blue'' and ``green''? Nature has no answer; languages do, and draw it differently. Russian has two basic color words for dark and light blue, siniy and goluboy. Native Russian speakers see their difference as you see blue versus red: ``two colors,'' not ``two shades of one color.'' Classical Chinese qīng comfortably spans blue and green: qīng sky is blue, qīng grass green, and qīng hair even black. Words are not printed on the rainbow. Each language takes its own knife to the continuous spectrum and cuts its own pieces.

Prototype theory is no laboratory toy; it is on duty around you every day. Think about ``good student.'' Your central prototype probably gets good grades, follows rules, and pleases teachers. But what about the top-scoring classmate who is late every day, or the one excelling in every subject who never raises a hand? Do they count? People hesitate, not because these students stand outside a particular score threshold, but because they are farther from the center. Many school injustices arise from mistaking ``far from the prototype'' for ``does not qualify.'' A student unlike the standard good student and a student who is not good are quite different things. Distinguishing them is prototype theory's first practical gift.

\textbf{Stop here: there is an especially easy mistake.} Many people hear about family resemblance and promptly conclude: ``So concepts are vague, and anything goes.'' That is wrong in three ways. First, a penguin is atypical but one hundred percent a bird. \textbf{Typicality comes in degrees; membership can be definite.} A circle with a center and margins is not no circle at all. Second, not every concept resists definition. ``Triangle'' has a perfect one: a closed figure enclosed by three line segments joined end to end. Meet it and you qualify; fail and you do not, with no exception in over two millennia. The same applies to ``odd number'' and ``square.'' Mathematical concepts are human-made precision tools; everyday concepts are wild hillsides. Different terrain: do not use one map for both routes. Third, societies can \textbf{construct boundaries} to handle vagueness. ``Adulthood'' is fuzzy in life, so law simply draws a line: eighteen years old, not one day short. The line is not ``natural,'' but it is useful. Family resemblance describes the real terrain of many everyday concepts. It is not a permit to ``say whatever you like.''

\section{Today: New Words Run Ahead, and Meaning Is Contested}
This old question plays out in your phone every day.

\textbf{The new-word production line is running at full speed.} Every few months another term appears: a small circle uses it first, then short videos and trending topics spread it to everyone. Finally, dictionary editors meet to discuss inclusion. If it enters, some celebrate that ``the times have been recorded'' while others protest that ``the dictionary has bowed to the internet.'' This is the tomato case replaying endlessly. Someone always holds a sign saying ``standards,'' someone else ``everyone uses it,'' while the word simply keeps moving.

\textbf{Old words quietly change address too.} Nèijuǎn, ``involution,'' began as an academic term: anthropologists used it for intensive cultivation with ever more input but no increase in output. Within a few years it was dragged into everyday language and expanded from ``internal competition without development'' to ``any effort that feels exhausting.'' Even memorizing two extra pages before an exam became involution. The cost of this expansion is that when a word holds everything, it is close to holding nothing. Other words suffer wounds on the battlefield. Yuàn once honorifically described a beautiful or distinguished woman, but repeated internet controversies dragged it through the mud. Now many people's first response to ``some kind of yuàn'' is sarcasm. Semantic drift is never merely linguistic; it records social feeling.

\textbf{Technical words are pulled into daily life.} ``Depression,'' medically an illness requiring diagnosis, now often means ``I'm in a bad mood today.'' ``OCD,'' originally a distressing mental disorder, becomes a playful way to say ``I like my books arranged neatly.'' ``Social anxiety'' shifts from a clinical concept to young people's self-mocking label. How should we judge this drift? Here our tools can help. On one hand, this is ordinary semantic broadening: dictionaries cannot and need not block it. On the other, it has real costs. When a diagnosis becomes a catchphrase, someone genuinely seeking help may hear ``I'm socially anxious too---who isn't?'' Diluted meaning quietly erases suffering. Usage shapes meaning, but users must bear the consequences when meaning thins. No institution can do this weighing for you; each speaker must do it.

\textbf{Nicknames are the politics of meaning nearest to you.} Nicknaming a classmate does exactly what this chapter's title warns against: stick a word onto a living person, then pretend that person has only one meaning. Once labeled, the whole person---who draws, tells jokes, falls at sports day and gets up again---is compressed into a sticky note. Understanding that words are not fixed labels gives you another reason for caution. Labels slide even on things; on people, they hurt.

\textbf{The newest player is artificial intelligence.} Large language models---the conversational AIs on your phone---learn language in precisely the way discussed here. Nobody makes them memorize dictionary definitions. They read enormous amounts of text, immersing every word in its uses. Having seen dǎ among millions of different neighbors, they too learn to glance around and identify which dǎ is meant. ``A word's meaning is its use in language'': AI resembles an engineering experiment on this philosophical assertion at unprecedented scale. But notice the gap. AI has usage's written traces, not the life behind them: it has never carried a water bucket, been hit, or played ball on a summer evening. It can therefore use words with extraordinary fluency yet sometimes produce statements elegant in usage and absurd in fact. The final centimeter between meaning and the real world comes from living, not reading.

\section{Over to You: Would You Want a ``Final Dictionary''?}
Return to the opening sticky note.

Suppose a ``language commission'' is established tomorrow with one task: publish the final dictionary. Every word retains exactly one official meaning, written into law. From then on, ``sugar-free'' must really mean no sugar, with fines for misuse. But new words also need approval before birth, old meanings are locked forever, poets may not apply ``spring'' to a person, and every new internet joke breaks the rules.

Would you support or oppose this dictionary? Give your answer and reasons. Before concluding, put all this chapter's weights on the scales.

You have the standardizers' case: stable meaning is the lifeblood of contracts, medicine instructions, and exams. Without oversight, fraudsters will be first to knead words into clay. Behind the French Academy and Iceland's invented tölva is a genuine feeling that ``language is our shared home.''

You have the usage side's case too: dictionaries always run behind language; English travels the world without police; nice drifted from ``foolish'' to ``pleasant'' without a hearing. And ``who regulates?'' is never a clean question: the pen defining meaning has historically belonged to the powerful.

You also know something deeper. Many everyday concepts---game, chair, vegetable---resist perfect definition, living through family resemblance and prototypes. Yet that does not mean anything goes: penguins remain birds, triangles remain definable, and established convention becomes a real rule.

So which is a better friend to ordinary people: stable meaning or freedom? Does a locked dictionary protect truth-tellers, or those holding the dictionary? Does a wholly free river of meaning lift boats, or fraudsters?

There is no standard answer. But notice this: every word you choose in answering helps shape what those words will mean tomorrow. The dictionary is neither the river's source nor its embankment; it is a chart drawn by those who come later. The river is you.
